% =====================================================================
%  CH 1-4 : Vision · Couche I et taxonomie · Couches et traçabilité · Carte
% =====================================================================
\section{Vision et objectifs}\label{sec:vision}

Le développement logiciel est entré dans l’ère \emph{agent-first} : une part croissante du
code est écrite, relue et modifiée par des modèles de langage plutôt que par des humains.
Cette bascule remet en question une hypothèse centrale de nos outils — que le code source
est optimisé pour la lecture humaine. TRIADE est la réponse disciplinée à ce réexamen :
plutôt que de supprimer les langages (une impossibilité de principe, puisqu’un modèle de
langage manipule des symboles textuels et que tout format binaire rendu lisible par un
agent constitue de facto un langage), TRIADE déplace le langage vers une fine couche de
spécification maintenue par l’agent, tandis que l’artefact livré est du code machine natif
pur. Le compilateur disparaît du pipeline de livraison : ce que l’agent écrit est ce qui
s’exécute — et rien d’autre.

Le nom encode l’architecture : trois couches indissociables — Intention (I), Binaire (B),
Carte (C) — unies par un verrou croisé cryptographique. La couche I contient le quoi et le
pourquoi : contrats typés, pré et post-conditions, invariants, cas d’échec, vecteurs de
vérification. La couche B contient le comment : des octets scellés, versionnés, à adresses
stables. La couche C est l’index naviguable qui lie les deux et sert d’unique interface de
développement à l’agent.

La v0.4.0 tire les enseignements de la première validation à l’échelle — le recodage
complet de RoomEq (démonstration v0.3) — et les convertit en normes : ce que la
démonstration a pratiqué (reverse-specs, pré-versions, partage de fixtures, mesure avant
optimisation) devient exigence ; ce que son analyse des défauts a révélé devient
taxonomie ; ce que son étude de cas établit et ne prouve pas reste explicitement hors
périmètre de preuve.

\subsection{Fonction objectif hiérarchisée — invariant de révision}\label{sec:objectif}

\begin{center}
\large\bfseries Sécurité $>$ Performance $>$ Optimisation $>$ Coût en tokens
\end{center}

\medskip

Arbitrage lexicographique : un objectif n’est jamais sacrifié au profit d’un objectif de
rang inférieur. La v0.4.0 consigne cette hiérarchie comme \textbf{invariant de révision}
(D19) : aucune révision future du protocole ne peut la modifier sans invalider
explicitement le protocole entier — pas l’ordre, pas les rangs, pas la non-compensation.

\begin{table}[htbp]
\centering\small
\caption{La fonction objectif et ses métriques de validation.}
\label{tab:objectif}
\begin{tabularx}{\textwidth}{@{}l Y Y@{}}
\toprule
\textbf{Objectif} & \textbf{Signification concrète} & \textbf{Métrique de validation} \\
\midrule
Sécurité & Sûreté mémoire validée par oracle ; intégrité carte $\leftrightarrow$ binaire ;
durcissement au chargement & 0 corruption détectée ; tout échec de verrou bloque la
session ; échec au lancement = arrêt propre \\
Performance & 100 \% code machine natif ; zéro abstraction intermédiaire &
Comparaison par bloc scellé face à une référence compilée \\
Optimisation & Descente à l’ISA sur hot paths identifiés — jamais sans mesure (R11) &
Gain mesuré et journalisé par bloc optimisé \\
Coût en tokens & Édition ciblée via la carte ; jamais de lecture intégrale du binaire &
Budgets par type d’opération (§\ref{sec:carte}) \\
\bottomrule
\end{tabularx}
\end{table}

\subsection{Périmètre et non-objectifs}\label{sec:perimetre}

Cinq non-objectifs. Premièrement, TRIADE ne prétend pas remplacer les langages dans
l’industrie : c’est une architecture de laboratoire. Deuxièmement, la furtivité face aux
EDR n’est pas recherchée : l’outil cible est un utilitaire de pentest pour tests autorisés
en laboratoire, appelant les API du système via les interfaces standards du profil.
Troisièmement, l’universalité est structurelle mais progressive : chaque profil doit être
instancié et validé (M6, M7) avant d’être déclaré supporté — le protocole ne promet pas la
portabilité, il la vérifie. Quatrièmement, aucune couche d’abstraction à l’exécution
(bytecode portable, machine virtuelle, traduction à la volée) n’est introduite : elle
trahirait la fonction objectif Performance. Cinquièmement, la couche de gouvernance
validée par la démonstration et le cœur natif du protocole (notation, verrou
cryptographique sur octets, patch à chaud) restent des niveaux de preuve distincts : les
jalons M0 à M7 demeurent la charge de validation du cœur natif.

\subsection{Ce que change la v0.4.0}\label{sec:nouveautes}

Dix-sept évolutions, toutes tracées dans la correspondance du §\ref{sec:correspondance} et
dans le journal des versions, toutes rattachées aux décisions D19 à D24 : normalisation de
la couche I, taxonomie des défauts, traçabilité I $\leftrightarrow$ C $\leftrightarrow$
B $\leftrightarrow$ tests, norme de fragmentation, localisation par intention, R10
renforcée, R11 nouvelle, suite de conformité par profil, parité inter-profils,
pré-versions, reverse-specs, journal d’optimisation, modèle de menace étendu, journal des
décisions normalisé, tableau de bord de conformité. Aucune règle existante n’est
affaiblie ; aucune couche n’est ajoutée à l’exécution ; la hiérarchie de la fonction
objectif est inchangée.

\section{La couche I : format normalisé des spécifications}\label{sec:coucheI}

La v0.3 décrivait la couche I comme « contrats légers » (D3), sans en fixer le gabarit —
la légèreté demeure, la structure devient normative (D19). Une spécification I n’est pas un
document libre : c’est un artefact diffable dont la structure conditionne la vérification.

\subsection{Identifiants stables}

Chaque spécification reçoit à sa création un identifiant \code{<DOM>-<NN>} : un domaine
parmi une liste figée par projet et un numéro séquentiel. L’identifiant est immuable et
jamais réutilisé ; un module supprimé laisse son identifiant vacant à jamais. C’est
l’équivalent exact, côté intention, de l’adresse stable côté binaire (R7) : les références
croisées — dépendances, carte, tests — ne cassent jamais. Le couple (identifiant,
révision) est la clé de traçabilité du §\ref{sec:traceabilite}.

\subsection{Gabarit obligatoire}

\begin{lstlisting}
<DOM>-<NN> . <nom court> . revision <R>
Intention : <une phrase — le pourquoi>
Contrat :
  pre  : P1..Pn    (preconditions numerotees, une phrase chacune)
  post : Q1..Qn    (postconditions numerotees)
  inv  : I1..In    (invariants)
Echecs :
  E1 : <condition> -> <sanction nommee : exception | arret propre | repli>
Vecteurs :
  V1 : <vecteur dore | propriete differentielle | test de rejet>
Fixtures    : <references aux fixtures contractuelles partagees>
Provenance  : ecrite-avant-code | reverse-spec de <binaire scelle @ tag>
Decisions   : <Dxx...>
\end{lstlisting}

Sept blocs, tous obligatoires — les échecs et vecteurs ne peuvent être vides que pour un
module sans effet observable, situation rare à justifier dans le champ Intention. Trois
règles de rédaction conditionnent la diffabilité exigée par R2 : \textbf{une exigence par
phrase numérotée} (le diff montre exactement quelle exigence change) ; \textbf{aucune
connaissance dupliquée} (renvoi par identifiant, sinon défaut T5) ; \textbf{chaque
pre/post/invariant relie au moins un vecteur} (une exigence invérifiable est une exigence
mal formulée). La provenance distingue les spécifications écrites avant le code — le cas
normal, exigé par R2 — des reverse-specs (§\ref{sec:reversespecs}).

\subsection{Révision et cycle de vie}

La révision s’incrémente à chaque changement du contrat. Une évolution qui rompt un
contrat existant ne modifie pas la spécification — elle en crée une nouvelle, l’ancienne
étant dépréciée avec renvoi : miroir exact du cycle de vie des blocs. Les appelants
dépendent du couple (identifiant, révision) ; la carte porte ce couple dans chaque entrée.

\subsection{Fixtures contractuelles partagées}

Quand plusieurs consommateurs se lient au même contrat externe — sortie d’un outil
système, format de fichier, API distante — les échantillons qui définissent ce contrat
appartiennent à la spécification, pas aux consommateurs. Le champ Fixtures référence un
jeu unique ; toute copie privée est un défaut de classe T5. Le mécanisme est né dans la
démonstration RoomEq : l’installateur scellé et l’application recodée partagent les
fixtures de parsing de pnputil — le défaut historique qui a survécu à trois versions
vivait précisément de la duplication.

\subsection{Taxonomie des défauts de contrat non écrit}\label{sec:taxonomie}

La classe de défauts que le protocole existe pour éliminer a un nom : \emph{contrat non
écrit}. La v0.4.0 en fixe la taxonomie (D19), née de l’analyse des défauts historiques de
la démonstration.

\begin{table}[htbp]
\centering\small
\caption{Taxonomie des défauts de contrat non écrit.}
\label{tab:taxonomie}
\begin{tabularx}{\textwidth}{@{}l Y Y@{}}
\toprule
\textbf{Classe} & \textbf{Définition} & \textbf{Verrou} \\
\midrule
T1 — valeur-seuil enfouie & Constante critique sans nom ni justification & constante
nommée + rationale dans la spec I \\
T2 — convention d’interface implicite & Structure ou ABI déclarée de mémoire & contrat
typé + invariant de taille testé \\
T3 — dépendance à l’hôte & Comportement dépendant de l’OS d’exécution & contrat
indépendant de l’hôte, union fixée \\
T4 — parsing de surface fragile & Analyse de libellés localisés ou volatils & fragments
indépendants des libellés + fixtures partagées \\
T5 — contrat dupliqué & Même connaissance codée en N exemplaires divergents & source
unique : spec I + fixtures \\
T6 — garantie non testée & Comportement critique sans vecteur de rejet & cas d’échec
nommé + test de rejet \\
\bottomrule
\end{tabularx}
\end{table}

Deux règles d’usage : tout défaut corrigé est classé dans cette taxonomie au journal des
décisions (la métrique « défauts ouverts par classe » du tableau de bord en dérive) ; la
taxonomie sert de liste de contrôle de rédaction — chaque bloc du gabarit verrouille au
moins une classe. Les occurrences de référence viennent de la démonstration : le seuil de
détection non justifié (T1+T6), la structure \code{SP\_DEVINFO\_DATA} de 424 octets au
lieu de 32 (T2, échec \code{0x6F8}), la sanitisation de noms dépendante de l’hôte (T3), le
parsing de pnputil sur libellés français (T4+T5, trois copies).

\section{Les trois couches : architecture et traçabilité}\label{sec:couches}

Le diff d’intention descend de I vers C, le verrou croisé relie C et B dans les deux sens,
et la couche B seule quitte le laboratoire. \textbf{Aucune couche n’est optionnelle} —
c’est une propriété normative : privé de I, le système perd son ancrage sémantique ; privé
de C, l’agent devient aveugle ; privé de B, il n’y a plus rien à exécuter. Toute
implémentation qui livrerait B sans I et C à jour viole R1 ; toute évolution qui
contournerait C pour éditer B directement viole R3 et la boucle.

\begin{table}[htbp]
\centering\small
\caption{Les trois couches.}
\label{tab:couches}
\begin{tabularx}{\textwidth}{@{}l l Y Y@{}}
\toprule
\textbf{Couche} & \textbf{Visibilité} & \textbf{Contenu} & \textbf{Rôle} \\
\midrule
I & Privée & Spécifications normalisées : contrats pré/post, invariants, cas d’échec,
vecteurs & Compression de l’intention \\
C & Privée, signée, chiffrée & Index des blocs, adresses, preuves, manifestes,
traçabilité & Unique interface de développement \\
B & Livrée & Code machine scellé, padding, checksums runtime & Exécution pure,
sémantiquement muette \\
\bottomrule
\end{tabularx}
\end{table}

La séparation croise une seconde dimension : la cible. La couche I est indépendante de la
cible — un contrat s’écrit une fois, ses types abstraits étant paramétrés par le profil.
La couche B est par profil. La couche C indexe l’ensemble : chaque entrée porte son
profil, et frag-bld organise les empreintes par profil. Quatre invariants structurels,
vérifiés à chaque cycle — leur violation est un refus de build : B n’est jamais nu (R1) ;
les adresses exposées sont stables (R7, par profil) ; la signature couvre la carte entière
(R9) ; B est auto-vérifiant au chargement (R10 renforcée).

\subsection{Traçabilité I $\leftrightarrow$ C $\leftrightarrow$ B $\leftrightarrow$
tests}\label{sec:traceabilite}

Chaque entrée de la carte référence quatre maillons : la spécification (identifiant +
révision), l’entrée d’index et le contrat embarqué, le(s) bloc(s) instancié(s) avec
empreinte(s) par profil, et les vecteurs exécutés avec le verdict du dernier cycle. Trois
propriétés sont contrôlées mécaniquement par la \textbf{commande d’audit de traçabilité} :
aucune spécification sans entrée de carte (l’exception visible étant la reverse-spec en
statut « prête ») ; aucun bloc scellé sans spécification référencée ; aucun vecteur
orphelin. Une chaîne incomplète est un refus de build — au même titre qu’un binaire nu.
La traçabilité est la condition de l’audit par des tiers : elle permet de répondre à
« quelle exigence ce bloc implémente-t-il, et quel test la défend ? » sans lecture du
binaire.

\section{La carte : fragmentation et budgets}\label{sec:carte}

\begin{sloppypar}
La carte est le cœur opérationnel du protocole : elle rend le binaire navigable, porte les
preuves et concentre la valeur intellectuelle. Elle est autodescriptive : un en-tête
normalisé consigne la version du protocole, la suite cryptographique — hachage, signature,
chiffrement, chacun dans une liste blanche (D15) — et les profils actifs. Sa spécification
détaillée fait l’objet de \code{FORMAT\_CARTE.md} — fragments, format d’entrée, verrou,
cycle de vie, export, re-cryptage —, mise à jour avec la présente révision : l’en-tête
porte \code{version=0.4.0} et chaque entrée de bloc porte le champ de traçabilité
\code{spec=<DOM>-<NN>:R}.
\end{sloppypar}

\paragraph{Norme de fragmentation.} La carte est fragmentée en quatre — frag-loc, frag-con,
frag-prf, frag-bld — et la v0.4.0 pose la norme : aucune opération de la boucle n’exige la
lecture de plus d’un fragment ; \textbf{aucune opération n’exige la lecture intégrale du
binaire} — le balayage est interdit (§\ref{sec:boucle}) ; l’audit complet de la carte est
une opération exceptionnelle, non budgétée, journalisée.

\paragraph{Budgets par type d’opération.}

\begin{table}[htbp]
\centering\small
\caption{Budgets en tokens par opération (inchangés depuis la v0.3, instrumentés dès M1).}
\label{tab:budgets}
\begin{tabularx}{\textwidth}{@{}Y Y l@{}}
\toprule
\textbf{Opération} & \textbf{Données lues} & \textbf{Budget} \\
\midrule
Localiser un bloc & frag-loc — une entrée & 40 tokens max \\
Éditer simplement un bloc & frag-loc + frag-con + notation & 3 000 tokens max \\
Ajouter un bloc & frag-loc + frag-con + seuils frag-prf & 6 000 tokens max \\
Instancier un profil nouveau & I partagée + frag-con + notation du profil & 6 000 tokens
max par bloc \\
Audit complet (rare) & carte entière déchiffrée & non budgété \\
\bottomrule
\end{tabularx}
\end{table}

Indicateurs de santé : précision de localisation (cible 95 \% dès M0), coût moyen d’une
édition simple (cible 5 000 tokens), taux de régression post-patch (cible zéro). Ces
budgets restent des cibles de conception jusqu’à leur instrumentation (M1) ; le tableau de
bord les affiche alors comme métriques mesurées.
